\chapter{Postać Frobeniusa i czynniki niezmiennicze}
\label{ch:lin19-frobenius}
\index{postać Frobeniusa@Postać Frobeniusa}
\index{postać Smitha@Postać Smitha}
\index{czynnik niezmienniczy@Czynnik niezmienniczy}
\index{moduł nad pierścieniem wielomianów@Moduł nad pierścieniem wielomianów}

\noindent\textbf{Cel rozdziału.}
Postać Jordana z Rozdziału~\ref{ch:lin17-jordan}
wymaga, by wielomian minimalny rozkładał się
nad rozpatrywanym ciałem na czynniki liniowe.
Przedstawimy postać kanoniczną dostępną nad
\emph{dowolnym} ciałem: postać Frobeniusa.
Jej bloki są macierzami towarzyszącymi wielomianów,
a podstawowym narzędziem jest postać Smitha
macierzy nad pierścieniem \(K[t]\).

W całym rozdziale \(K\) jest dowolnym ciałem,
\(V\) przestrzenią wektorową nad \(K\)
o wymiarze \(n<\infty\), zaś
\(T\in\operatorname{End}_K(V)\).
Dopuszczamy \(V=\{0\}\), z pustą bazą,
pustą macierzą i \(\chi_T=m_T=1\).

\section{Endomorfizm jako moduł nad \(K[t]\)}
\label{sec:lin19-module}

\begin{definicja}[Moduł operatora]
\label{def:lin19-module}
Na przestrzeni \(V\) określamy działanie
pierścienia wielomianów \(K[t]\) wzorem
\[
 p\cdot v=p(T)v \qquad (p\in K[t],\ v\in V).
\]
Moduł otrzymany z tego działania
oznaczamy przez \(V_T\).
\end{definicja}

\begin{stwierdzenie}[Podobieństwo i izomorfizm modułów]
\label{prop:lin19-module-similarity}
Jeżeli \(T_i\in\operatorname{End}_K(V_i)\),
\(i=1,2\), to \(V_{T_1}\cong V_{T_2}\)
jako moduły nad \(K[t]\) wtedy i tylko wtedy,
gdy istnieje izomorfizm \(K\)-liniowy
\(U\colon V_1\to V_2\) spełniający
\(UT_1=T_2U\).
W szczególności dla macierzy kwadratowych
jest to równoważne podobieństwu.
\end{stwierdzenie}
\begin{proof}
Dla wielomianów \(p,q\in K[t]\)
mamy \(1(T)=I_V\) oraz
\((p+q)(T)=p(T)+q(T)\).
Ponadto \(p(T)q(T)=(pq)(T)\):
po zapisaniu \(p=\sum_i a_it^i\)
i \(q=\sum_jb_jt^j\)
obie strony są sumą
\(\sum_{i,j}a_ib_jT^{i+j}\).
Zatem działanie spełnia aksjomaty modułu.

Izomorfizm modułów jest \(K\)-liniowy,
ponieważ \(K\subset K[t]\) jako wielomiany stałe.
Zachowując działanie wielomianu \(t\),
otrzymujemy \(UT_1=T_2U\).
Odwrotnie, z tej równości indukcyjnie
wynika \(UT_1^k=T_2^kU\) dla \(k\ge0\),
a po liniowym zsumowaniu
\(Up(T_1)=p(T_2)U\) dla każdego wielomianu.
To właśnie warunek liniowości nad \(K[t]\).
\end{proof}

\begin{twierdzenie}[Prezentacja modułu operatora]
\label{thm:lin19-presentation}
Niech \(A\in\operatorname{Mat}_{n\times n}(K)\)
będzie macierzą \(T\) w pewnej bazie \(V\).
Wówczas istnieje izomorfizm modułów
\[
 V_T\cong K[t]^n/(tI_n-A)K[t]^n.
\]
Dla \(n=0\) obie strony są zerowe.
\end{twierdzenie}
\begin{proof}
Identyfikujemy \(V\) z \(K^n\).
Dla wektora wielomianowego
\(f(t)=\sum_{k=0}^{m}t^kv_k\),
gdzie \(v_k\in K^n\), definiujemy
\[
 \varepsilon(f)=\sum_{k=0}^{m}A^kv_k.
\]
Jest to odwzorowanie surjektywne,
bo stały wektor wielomianowy \(v\)
ma obraz \(v\).
Spełnia \(\varepsilon(tf)=A\varepsilon(f)\)
i jest \(K\)-liniowe, zatem
\(\varepsilon(pf)=p(A)\varepsilon(f)\)
dla wszystkich \(p\in K[t]\).
Oznacza to liniowość nad \(K[t]\).

Jeżeli \(g\in K[t]^n\), to
\[
 \varepsilon((tI_n-A)g)
 =A\varepsilon(g)-A\varepsilon(g)=0.
\]
Stąd \((tI_n-A)K[t]^n\subset\ker\varepsilon\).
Dla odwrotnego zawierania zauważmy, że
dla każdego \(k\ge1\) zachodzi
\[
 t^kI_n-A^k
  =(tI_n-A)\sum_{j=0}^{k-1}t^{k-1-j}A^j.
\]
Po wymnożeniu składniki pośrednie
skracają się parami; pozostają
jedynie \(t^kI_n\) oraz \(-A^k\).
Mnożąc przez \(v_k\) i sumując
po \(k=1,\ldots,m\), otrzymujemy
\[
 f(t)-\varepsilon(f)
 =(tI_n-A)\sum_{k=1}^{m}\sum_{j=0}^{k-1}
          t^{k-1-j}A^jv_k.
\]
Jeśli \(\varepsilon(f)=0\),
to \(f\in(tI_n-A)K[t]^n\).
A zatem \(\ker\varepsilon\)
jest dokładnie tym podmodułem.
Izomorfizm \(K[t]^n/\ker\varepsilon\to V_T\)
przyporządkowujący klasie \(f\)
wartość \(\varepsilon(f)\) jest
dobrze określony, różnowartościowy,
surjektywny i liniowy nad \(K[t]\).
\end{proof}

\section{Postać Smitha nad pierścieniem \(K[t]\)}
\label{sec:lin19-smith}

\begin{definicja}[Operacje elementarne nad \(K[t]\)]
\label{def:lin19-elementary}
Na wierszach lub kolumnach macierzy
o elementach w \(K[t]\) dopuszczamy:
zamianę dwóch wierszy lub kolumn,
pomnożenie jednego z nich
przez niezerowy skalar z \(K\)
oraz dodanie do jednego wiersza
lub kolumny wielomianowej wielokrotności innego.
Odpowiadające im macierze są odwracalne
nad \(K[t]\), ponieważ każda operacja
ma odwrotność tego samego rodzaju.
\end{definicja}

\begin{twierdzenie}[Postać Smitha i jej jednoznaczność]
\label{thm:lin19-smith}
Niech \(M\in\operatorname{Mat}_{n\times n}(K[t])\)
oraz niech \(\det M\ne0\).
Istnieją macierze odwracalne
\(U,W\in\operatorname{GL}_n(K[t])\)
i moniczne wielomiany
\(d_1,\ldots,d_n\in K[t]\), takie że
\[
 UMW=\operatorname{diag}(d_1,\ldots,d_n),
 \qquad d_1\mid d_2\mid\cdots\mid d_n.
\]
Wielomiany \(d_i\) są jednoznacznie
wyznaczone przez \(M\).
Dla \(n=0\) ciąg wielomianów jest pusty.
\end{twierdzenie}
\begin{proof}
\emph{Istnienie.}
Przy \(n=0\) wynik jest zgodny
z przyjętymi konwencjami.
Niech \(n\ge1\).
Ponieważ \(\det M\ne0\),
macierz ma niezerowy element.
Za pomocą zamian wierszy i kolumn
przenieśmy taki element
na miejsce \((1,1)\) i oznaczmy go \(d\).

Dzielmy kolejne elementy
pierwszej kolumny przez \(d\).
Gdy element w pewnym wierszu
nie dzieli się przez \(d\),
odejmujemy wielomianową wielokrotność
pierwszego wiersza, pozostawiając
niezerową resztę w pierwszej kolumnie,
a następnie zamieniamy te dwa wiersze.
Stopień niezerowego elementu na miejscu
\((1,1)\) ściśle maleje.
Tak samo postępujemy z pierwszym
wierszem, używając operacji na kolumnach.
Powtarzamy dzielenie i zamiany,
aż \(d\) dzieli wszystkie elementy
pierwszego wiersza i kolumny.
Zakończenie wynika z faktu, że
stopień wielomianu niezerowego
jest nieujemną liczbą całkowitą
i każda nieudana próba dzielenia
obniża stopień elementu wyróżnionego.
Odejmując odpowiednie wielokrotności,
doprowadzamy macierz do postaci
\[
 \begin{pmatrix}d&0\\0&B\end{pmatrix}.
\]

Musimy jeszcze zapewnić, że \(d\)
dzieli wszystkie elementy \(B\).
Jeśli tak nie jest, wybierzmy element
\(b\) w bloku \(B\), który nie jest
podzielny przez \(d\).
Dodajemy do pierwszego wiersza
wiersz zawierający \(b\).
Na miejscu \((1,1)\) pozostaje \(d\),
natomiast w pierwszym wierszu
pojawia się \(b\).
Dzielenie \(b\) przez \(d\),
odjęcie wielokrotności pierwszej kolumny
i zamiana kolumn pozwalają otrzymać
nowy niezerowy element w pozycji \((1,1)\)
o stopniu mniejszym od \(\deg d\).
Wznawiamy opisany algorytm.
Ponieważ kolejne takie powroty
zmniejszają stopień wyróżnionego elementu,
ostatecznie uzyskujemy
\(\operatorname{diag}(d,B)\)
z \(d\) dzielącym wszystkie elementy \(B\).

Wyznacznik operowanej macierzy
jest niezerowy: operacje elementarne
mnożą go jedynie przez odwracalny
element \(K[t]\), czyli przez
niezerową stałą.
Wobec tego \(\det B\ne0\)
dla \(n\ge2\).
Indukcyjnie stosujemy konstrukcję
do macierzy \(B\).
Wszystkie jej elementy są wielokrotnościami
\(d\), a operacje elementarne zachowują
tę własność. Dlatego uzyskane
elementy diagonalne dzielą się przez \(d\).
Mnożąc każdy wiersz przez odwrotność
współczynnika wiodącego, normalizujemy
wielomiany diagonalne do postaci monicznej.
Dostajemy \(d_1\mid\cdots\mid d_n\).
Wszystkie operacje odpowiadają
mnożeniu z lewej i prawej
przez macierze odwracalne nad \(K[t]\).

\emph{Jednoznaczność.}
Dla \(1\le k\le n\) niech \(\Delta_k(M)\)
oznacza moniczny największy wspólny dzielnik
wszystkich minorów stopnia \(k\) macierzy \(M\);
przyjmujemy \(\Delta_0(M)=1\).
Ponieważ \(\det M\ne0\),
rodzina minorów każdego stopnia \(k\)
zawiera element niezerowy:
niezerowy wyznacznik wymusza
rangę \(n\) nad ciałem ułamków \(K(t)\),
a więc istnieje niezerowy minor
każdego stopnia \(k\).

Dodanie wielomianowej wielokrotności
jednego wiersza do drugiego
zmienia każdy minor w kombinację liniową
minory tego samego stopnia,
o współczynnikach w \(K[t]\);
analogicznie działa operacja na kolumnach.
Zamiana i mnożenie przez niezerową stałą
również zachowują ideał generowany przez minory.
Ponieważ operacje są odwracalne,
ideały generowane przez minory przed
i po każdej operacji są dokładnie równe.
W konsekwencji wszystkie \(\Delta_k\)
są niezmiennikami operacji elementarnych.

Dla macierzy diagonalnej z warunkami
\(d_1\mid\cdots\mid d_n\)
każdy niezerowy minor stopnia \(k\)
jest iloczynem \(k\) wybranych
elementów diagonalnych.
Ponieważ \(d_i\mid d_j\) dla \(i\le j\),
iloczyn \(d_1\cdots d_k\)
dzieli wszystkie takie minory.
Sam jest jednym z minorów,
zatem
\[
 \Delta_k(M)=d_1\cdots d_k,\qquad
 d_k=\frac{\Delta_k(M)}{\Delta_{k-1}(M)}.
\]
Wzór ten jednoznacznie określa
moniczne wielomiany \(d_k\).
\end{proof}

\begin{przyklad}[Dwa czynniki macierzy wielomianowej]
\label{prz:lin19-smith-jordan}
Niech
\[
 A=\begin{pmatrix}0&1\\0&0\end{pmatrix},
 \qquad
 M(t)=tI_2-A=\begin{pmatrix}t&-1\\0&t\end{pmatrix}.
\]
Największy wspólny dzielnik
elementów \(M\) jest równy \(1\),
gdyż występuje element \(-1\).
Zatem \(\Delta_1(M)=1\).
Jedyny minor stopnia \(2\)
to \(\det M=t^2\);
otrzymujemy \(\Delta_2(M)=t^2\).
Postać Smitha ma więc
elementy diagonalne \(1,t^2\).
\end{przyklad}

